\section{One lease total can conceal different rights}
\label{sec:lease-allocation-family}

The preceding lifecycle starts with a supplied asset of 30 for right B.
Suppose that amount is unknown and only the combined right-of-use asset of 120
is given. Even if the two rights and their liabilities of 100 and 50 are known,
the total asset does not identify its allocation. Assigning 20 to B leaves 100
for A; assigning 30 to B leaves 90 for A. Both allocations produce exactly the
same opening signed account balances. Terminating B gives a gain of 30 in the
first case and 20 in the second.

Let $x$ be B's opening right-of-use asset, in the same exact USD units as the
previous example. A's asset is then $120-x$. Nonnegative asset amounts imply
\[
  0\leq x\leq120,\qquad R_A=120-x,\qquad R_B=x.
\]
This is a family of supplied hypothetical openings. Exactly two identified
rights, their liability allocation, the fixed 10 percent rate and the original
method findings remain premises. Actual aggregate disclosures need not establish
those premises. The family assigns no probabilities to its allocations.

\subsection{Find where the supplied lifecycle can complete}

Keep the original four events, including the annual payments and depreciation
of 81. Only the unknown opening asset allocation and B's corresponding traced
removal vary. Terminating B removes liability 50 and asset $x$, so its gain is
$50-x$. A's payment of 121 gives liability 110 and adds 10 to its asset. After
the first event, total ROU is therefore $130-x$.

The second reduction removes asset 10 and liability 11 from A. Remeasuring its
remaining payment of 99 gives liability 90, which is 9 below the intermediate
liability of 99. Its asset becomes
\[
  R_A^{(2)}=(130-x)-10+(90-99)=111-x.
\]
The original method requires a nonnegative remaining asset. For $x>111$, this
second reduction is outside its supported treatment. That refusal does not make
the opening allocation impossible; it identifies a method limitation under the
supplied event. The first completed event remains available.

For $x\leq111$, the payment of 99 settles liability 90, records interest expense
9 and leaves company cash 101. The final event then proposes depreciation of
81. It can be applied only when that expense does not exceed the remaining
asset:
\[
  81\leq111-x\quad\Longleftrightarrow\quad x\leq30.
\]
Thus the complete supplied sequence works on $[0,30]$. For $30<x\leq111$, it
reaches payment but rejects the final depreciation. Replacing 81 with a smaller
amount would change the supplied scenario; the reference keeps the refusal and
its successful prefix instead.

The three regions partition the original allocation interval:
\begin{center}
\begin{tabular}{lll}
\toprule
Allocation of B's asset & Outcome & Completed events \\
\midrule
$[0,30]$ & Complete supplied lifecycle & Four \\
$(30,111]$ & Depreciation amount rejected & Three \\
$(111,120]$ & Second reduction unsupported & One \\
\bottomrule
\end{tabular}
\end{center}
The endpoints matter. At $x=30$, depreciation removes the remaining asset
exactly. At $x=111$, the second reduction leaves zero ROU, so it succeeds, but
the later depreciation of 81 does not. Each failed event leaves the last
successful state unchanged. A response containing only the first region would
lose valid declared openings and conceal why their subsequent paths stopped.

\subsection{Keep financial outcomes on the same allocation}

Within the complete region, depreciation leaves asset $30-x$. Pretax profit
combines the first gain $50-x$, the second gain 1, interest 9 and depreciation
81. Write $R_{\rm end}$ for final ROU and $\Pi$ for cumulative pretax profit:
\[
  R_{\rm end}=30-x,\qquad
  \Pi=(50-x)+1-9-81=-39-x,\qquad 0\leq x\leq30.
\]
The possible asset amounts range from 0 to 30 and profits from $-69$ to $-39$.
Both depend on the same opening allocation. Subtracting the two expressions
eliminates $x$:
\[
  \Pi-R_{\rm end}=(-39-x)-(30-x)=-69.
\]
For example, $x=20$ produces asset 10 and profit $-59$; $x=30$ produces asset
zero and profit $-69$. Asset 30 and profit $-69$ are individually feasible
marginal endpoints, but they cannot occur together. The first requires $x=0$;
the second requires $x=30$.

A hypothetical additional observation of final ROU equal to 10 would select
$x=20$ and profit $-59$, conditional on completion of the supplied lifecycle.
The original unconditioned family still contains its rejected and unsupported
regions. The query records the additional observation as a supplied premise;
it does not present that number as a recovered issuer fact.

The interrupted paths also retain dependent financial amounts. In the
rejected-depreciation region, the last successful state has ROU $111-x$, profit
$42-x$, cash 101 and zero lease liability. In the unsupported-reduction region,
the first event leaves ROU $130-x$, profit $50-x$, cash 200 and signed lease
liability $-110$. These are prefix states, not completed-year statements. Keeping
their dates and failure positions prevents a partially executed path from being
reported as a complete result.

\subsection{Represent the family without dropping interrupted paths}

\path{LeaseAllocationFamily/1} binds the original source lock, a finite
allocation interval within $[0,120]$, explicit endpoint inclusion, an allocation
premise, its source reference, a complete-membership finding and a
\path{max_cells} budget. Amounts are canonical rational strings. The reference
intersects the declared interval with all three regions. Empty intersections
contribute no cell; every nonempty intersection remains visible. Missing
membership or a missing allocation premise returns unresolved.

For each checked cell, the response retains its domain, original outcome
category, completed-prefix length, failure index and reason, affine financial
expressions and a complete native replay. The replay contains opening state,
events, proposals and the last successful state. It checks a representative
allocation against the unchanged original methods. The inspected affine
derivation establishes the within-cell relationship; a few replayed allocations
alone would not prove a bound over a continuous interval.

The cell budget limits those native replays. If it is exhausted, the remaining
domains and categories stay present as unchecked cells. The response reports
partial search, supplies no complete-family bounds and makes no assertion that
all allocations complete. Once every retained cell is checked, any final-outcome
bounds are explicitly conditional on complete execution. A family consisting
entirely of interrupted paths has no such bounds.

\path{LeaseAllocationQuery/1} requires the explicit condition
\path{complete_execution}. Its outcome constraints intersect one shared
allocation interval, preserving dependence between asset, profit, cash and
liability. The response retains the unconditioned family alongside the selected
interval and its conditional ranges. An impossible combination has an empty
selection and no bounds. A separate \path{max_replays} budget controls the
query witness: with zero budget, a nonempty candidate interval remains visible
but its selection and bounds are not reported as verified.

The source locks and full recomputation reject changed methods or forged saved
results. Source admission, accounting admission, method eligibility, actual
history identification and production commit remain false. The exact original
dates, rates, schedules and depreciation are restrictions of this reference;
unknown rates or liabilities, arbitrary membership and broader lease valuation
require additional methods and cannot be handled by renaming this profile.

The derivation and complete saved requests and responses are retained under
\path{data/accounting-master/}:
\begin{quote}\small
\path{lease-allocation-family-audit-20261024}\\
\path{lease-allocation-family-20261024}
\end{quote}
They show how aggregate-compatible openings can remain distinct through a
supported accounting calculation, and how unsuccessful paths remain part of
the answer. Connecting that family to actual disclosed rights, measurements
and authority is still a separate source-reconstruction task.
